\documentclass[11pt,a4paper]{article}
\usepackage[margin=25mm]{geometry}
\usepackage{amsmath,amssymb,amsthm}
\newtheorem{theorem}{Theorem}
\title{Rank spectra for the Coxeter source}
\author{Human Compression project: F006 recovery}
\date{Assembly for review}
\begin{document}\maketitle
This manuscript joins three complete components concerning the principal
results of Chen and Rowley, \emph{Abstract Regular Polytopes of Finite
Irreducible Coxeter Groups}, Electronic Journal of Combinatorics 32(4)
(2025), P4.18. We use the fixed thirteen-page published edition.
The source's definition permits a string diagram to be disconnected;
its discussion of $B_n$ explicitly uses that convention. All statements
below retain it unless connectedness is specified.

\begin{theorem}
For $n\ge5$, the group $W(D_n)$ admits string C-markings of every
rank from three to $n$ when $n$ is odd, and from three to $n-1$
when $n$ is even. Those upper endpoints are its maximum ranks.
For the exceptional groups, the maximum ranks are
\[
\begin{array}{c|rrrrrr}
W&H_3&H_4&F_4&E_6&E_7&E_8\\\hline
r_{\max}(W)&4&4&4&5&7&7\\
\text{maximum connected rank}&3&4&4&5&6&7.
\end{array}
\]
For $I_2(m)$, $m\ge3$, the maximum rank is three when
$m\equiv2\pmod4$ and two otherwise. Its connected maximum
is always two. The six connected exceptional censuses agree
with every entry of the source's Table~1. With disconnected
strings included, $H_3$ has three additional classes at rank
four, and $E_7$ has $462,276,72,18$ additional classes at ranks
four, five, six and seven, respectively, all modulo isomorphism
and duality. None of these additional classes is self-dual.
\end{theorem}
\begin{proof}
The acquired \emph{Even sign changes and the rank of a string} proves the
upper bounds for $W(D_n)$ through its even sign module and the
independent generating families in the symmetric quotient. Its
companion \emph{Constructing the rank interval with sign changes}
proves the displayed constructions and their full intersection
property. In odd degree, the acquired contraction theorem gives
every intermediate rank down to three; in even degree, the
explicit families give rank three and each rank from four to
$n-1$. These statements apply to all the specified degrees, with
no finite-degree extrapolation.

The acquired \emph{Disconnected strings in two Coxeter families}
proves the uniform dihedral assertion, the explicit $H_3$
rank-four witness, and the complete small census. Theorem~1
of \emph{The exceptional Coxeter censuses}, with its full
enumeration proofs and retained executions, gives the other five
groups. Their final nonzero rows give both kinds of maximum.
The separate direct subset-intersection verification also supplies
an explicit disconnected rank-seven $E_7$ witness. Subtracting
the connected $E_7$ rows from its complete rows gives
$1987-1525=462$, $741-465=276$, $167-95=72$ and $18-0=18$.
The self-dual columns agree in each subtraction. The actual
comparison with the source's six printed rows is retained,
including all zero entries and the total columns.
\end{proof}

\paragraph{The exact recovery contract.}
The first assertion recovers source Theorems~1 and~2, including
both the upper bounds and all intermediate ranks. The remaining
assertions recover Theorem~3 and Table~1 with their disconnected
scope corrected in precisely $I_2(m)$, $H_3$ and $E_7$.
The separate $D_n$ source map records every intermediate
numbered statement and retained input: in particular, its new
kernel calculation corrects a displayed intermediate intersection
in Lemma~19 while preserving that lemma's rank-three conclusion.
These are corrections to fixed published assertions, not answers
to open conjectures.

The source's constructions, conclusions and connected counts
remain attributed to Chen and Rowley. The new arguments are the
explicit even-module and kernel derivations, the treatment of
disconnected strings, and the complete new finite executions.
Classical Coxeter identifications and orders, the cited
symmetric-group generation results, the string end criterion,
and the acquired rank-contraction theorem remain stated inputs.
The enumeration follows the established Hartley--Hulpke and
Leemans--Mixer method with the exact implementations and their
proofs retained. No result about the cited external literature
as a whole, historical priority, or other editions is asserted.

The proofs depend on the specified exact Python and GAP
computations where explicitly stated. Those programs are
computational premises and checks; they are not complete Lean
formalizations. This assembly introduces no compression
measurement and does not turn the source's corrected assertions
into claimed literal recoveries of their false unrestricted form.
\end{document}
